\section{稳定区域与刚性方程}
\label{sec:08-Numerical-Analysis/06-ODE/03}

一个化学反应器的仿真跑了一夜没跑完。日志显示求解器接受的步长长期停在 \(2\times10^{-6}\)，而画出来的浓度曲线在那段时间里平得像一条直线。把容差从 \(10^{-8}\) 放宽到 \(10^{-4}\)，步长纹丝不动。

把容差放宽而步长不动，说明限制步长的不是精度。上一节的所有工具在这里全部失效：误差估计明明合格，步长却涨不上去，因为一旦涨上去，数值解会开始做真解从未做过的事。

\subsection{稳定性和精度量的不是同一件事}

判断这件事只需要一个方程。取

\[
y'=\lambda y,\qquad \operatorname{Re}\lambda<0,
\]

真解 \(e^{\lambda t}\) 衰减到零。绝大多数一步法作用在它上面都会化成一次乘法，

\[
y_{n+1}=R(z)\,y_n,\qquad z=h\lambda,\qquad y_n=R(z)^n y_0 .
\]

\(R\) 叫方法的稳定函数。真解衰减而数值解不衰减，当且仅当 \(\abs{R(z)}>1\)。于是把

\[
\mathcal S=\set{z\in\C:\abs{R(z)}\le1}
\]

叫作这个方法的绝对稳定区域。它是方法的属性，与具体问题无关；具体问题只提供 \(\lambda\)，你只提供 \(h\)，两者相乘落在 \(\mathcal S\) 里面还是外面，决定数值解是跟着衰减还是掉头长大。

显式 Euler 的 \(R(z)=1+z\)，稳定区域是以 \(-1\) 为心、半径 \(1\) 的圆盘。后向 Euler 的 \(R(z)=1/(1-z)\)，\(\abs{1-z}\ge1\) 的区域包含整个左半平面。两者的差别不是量的差别。

\begin{center}
\begin{tikzpicture}[x=0.82cm,y=0.82cm,font=\small]
  \begin{scope}
    \draw[fill=blue!16,draw=blue!60] (-1,0) circle (1);
    \draw[->,black!55] (-3.5,0) -- (1.5,0) node[right,font=\footnotesize] {\(\operatorname{Re}z\)};
    \draw[->,black!55] (0,-1.9) -- (0,1.9) node[above,font=\footnotesize] {\(\operatorname{Im}z\)};
    \draw[black!55] (-2,-0.12) -- (-2,0.12);
    \node[below,font=\footnotesize,black!70] at (-2,-0.14) {\(-2\)};
    \fill[blue!75] (-0.7,0.4) circle (1.7pt);
    \fill[red!75] (-2.9,0.7) circle (1.7pt);
    \node[above,font=\footnotesize] at (-1,2.15) {显式 Euler：\(R(z)=1+z\)};
    \node[font=\footnotesize,black!70] at (-1,-2.6) {圆盘之外的 \(h\lambda\) 一律发散};
  \end{scope}
  \begin{scope}[shift={(7.2,0)}]
    \fill[blue!16] (-3.5,-1.9) rectangle (0,1.9);
    \draw[blue!60] (0,-1.9) -- (0,1.9);
    \draw[fill=red!14,draw=red!55] (1,0) circle (1);
    \draw[->,black!55] (-3.5,0) -- (2.4,0) node[right,font=\footnotesize] {\(\operatorname{Re}z\)};
    \draw[->,black!55] (0,-1.9) -- (0,1.9) node[above,font=\footnotesize] {\(\operatorname{Im}z\)};
    \fill[blue!75] (-0.7,0.4) circle (1.7pt);
    \fill[blue!75] (-2.9,0.7) circle (1.7pt);
    \node[above,font=\footnotesize] at (-0.6,2.15) {后向 Euler：\(R(z)=1/(1-z)\)};
    \node[font=\footnotesize,black!70] at (-0.6,-2.6) {整个左半平面都在区域内};
  \end{scope}
\end{tikzpicture}

\emph{同样两个 \(h\lambda\)，左图里远的那个落在圆盘之外，数值解掉头增长；右图里它仍在阴影内。隐式方法的不稳定区域反而缩到右半平面的一个小圆盘上——那里的真解本来就在增长。}
\end{center}

图里两个点位置没变，动的只是方法。左图的红点表示：真解在指数衰减，数值解在指数增长，误差不是大了几倍，是符号和趋势都反了。回到上一节末尾那个 \(y'=-y\)、\(h>2\) 的例子，那正是 \(z=-h\) 越过了圆盘左端的 \(-2\)。

要留意图中蓝色区域并不承诺精度。落在区域内只说明不会出现这种虚假增长；用后向 Euler 迈一个巨大的步长，可以极其稳定地把一条有起伏的真解抹成一条平滑的假曲线。稳定与准确是两把尺子，通过一把不代表通过另一把。

\subsection{步长被一个你并不关心的分量绑架}

单个 \(\lambda\) 还看不出问题的严重性。把两个模态放在一起：系统的 Jacobian 有特征值 \(\lambda_1=-1\) 和 \(\lambda_2=-10^{6}\)。

快模态在 \(t\approx10^{-5}\) 之内就衰减到机器精度以下，此后的解完全由慢模态主导，特征时间尺度是 \(1\)。想画出这条慢曲线，\(h=10^{-2}\) 绰绰有余。可显式 Euler 要保持稳定，必须让两个 \(h\lambda_i\) 同时落在圆盘里，于是

\[
h\le\frac{2}{\abs{\lambda_2}}=2\times10^{-6}.
\]

约束来自 \(\lambda_2\)。而 \(\lambda_2\) 对应的那个分量，在积分开始后的第一微秒就已经死掉了，它对答案的贡献是零。步长被一个数值上已经不存在的东西按住，整整五个数量级。要在 \([0,10]\) 上积完，得走五百万步；换成隐式方法，几千步就够。

这就是刚性的操作性含义：显式方法的稳定步长远小于解的慢变部分所要求的精度步长，而且差距大到改变计算的可行性。它不是``解变化得快''——一条剧烈振荡的解可以完全不刚性，因为小步长本来就是精度所需；也不是``导数数值很大''。刚性是问题、时间区间与方法三者共同的性质，大负实部特征值与时间尺度分离是最常见的信号，但不存在一个适用于所有非线性非正规系统的标量定义。

判断求解器是不是卡在这里，有一个很利落的实验：把容差放宽两个数量级。若步长跟着变大，它受精度约束；若步长几乎不动，它受稳定性约束，继续调容差是白费力气，该换方法类别了。

刚性在训练里有一个几乎逐字对应的版本。梯度下降作用在二次近似上是 \(w_{k+1}=(I-\eta H)w_k\)，收敛要求所有 \(\abs{1-\eta\mu_i}<1\)，也就是 \(\eta<2/\lambda_{\max}(H)\)——形式与上面那条 \(h\le2/\abs{\lambda_2}\) 完全一样。学习率被最陡的那个方向卡死，而最平坦的方向按 \(1-\eta\mu_{\min}\) 的速度缓慢移动，收敛所需的步数正比于条件数 \(\lambda_{\max}/\mu_{\min}\)。参数尺度差异极大的模型就是一个刚性系统，\hyperref[sec:08-Numerical-Analysis/01-Floating-Point-and-Error/03]{``条件数衡量问题本身的敏感性''}和\hyperref[sec:09-Convex-Optimization/06-Optimization-Algorithms/01]{``梯度下降真正难在步长''}从两个方向讲的是同一件事。预条件、归一化层和自适应优化器所做的，相当于这里换用隐式方法。

\subsection{隐式方法用一次解方程买下整个左半平面}

后向 Euler 每步要解

\[
y_{n+1}=y_n+h f(t_{n+1},y_{n+1}),
\]

未知量在等号两边。对线性问题这是一次线性求解，对非线性问题是一轮 Newton 迭代：反复解 \(J\Delta=-F\) 再更新 \(y\)，其中 \(J=I-h\,\partial f/\partial y\)。步长越大，这个矩阵越远离单位阵，迭代越吃力，但稳定性上的收益远远盖过这点成本。

梯形法取两端斜率的平均，

\[
y_{n+1}=y_n+\frac h2\bigl[f(t_n,y_n)+f(t_{n+1},y_{n+1})\bigr],
\qquad R(z)=\frac{1+z/2}{1-z/2},
\]

也覆盖整个左半平面，而且是二阶的。两者都称为 A-stable。差别出现在 \(z\to-\infty\) 的极限上：后向 Euler 的 \(R(z)\to0\)，把极快模态一步碾平，称为 L-stable；梯形法的 \(R(z)\to-1\)，快模态的幅值不衰减，只在相邻步之间来回变号，留下一串肉眼可见的锯齿。A-stable 不蕴含 L-stable，这条区别在强刚性问题上会以``解看起来在抖''的形式找上门。

生产级的刚性求解器（BDF 族、Radau 隐式 RK、SDIRK）都在这条线上，成本集中在三处：Jacobian 的构造或近似、稀疏线性系统的求解、以及矩阵分解的复用策略。步长与 Jacobian 变化不大时复用已有分解，是这类代码里最关键的一项工程。大型系统里线性求解本身要交给 Krylov 方法配预条件，见\hyperref[sec:08-Numerical-Analysis/05-Numerical-Linear-Algebra/04]{``大型稀疏系统与 Krylov 迭代''}。

隐式不等于免费。非线性迭代可能不收敛；Jacobian 近似得太糙会把方法的实际阶数拖下来；内层容差过松会污染时间积分，过紧则把时间浪费在解一个本来就只需近似的方程上。选求解器时该比的是总函数求值次数、Jacobian 次数、线性迭代次数和实际墙钟时间，只看步数会得出反过来的结论。

\subsection{稳定图能诊断，不能担保}

最后收一下边界。\(R(z)\) 的分析建立在标量线性测试方程上，推广到方程组时默认 Jacobian 可对角化且各模态解耦。对非正规矩阵，即使全部特征值都在左半平面，解也可能先经历一段显著的暂态增长，特征值给出的只是长期行为；此时该看的是伪谱或数值域，而不是那张圆盘图。非线性系统的 Jacobian 随轨迹变化，某一点的局部谱不构成全局的证明。

即便如此，稳定函数仍然是遇到``求解器为什么这么慢''时值得问的第一个问题，因为它区分开两种完全不同的处方：精度不够就调容差或换高阶方法，稳定性不够就得换到隐式方法那一类去。答错这道题，优化方向会整个反掉。

到这里，沿时间推进的那条线索基本走完了。下一节换一类问题：约束不在起点，而是分散在区间两端，逐步推进连出发的方向都定不下来。
